\section{Related work}
\label{sec:related}

\subsection{Pruning}

\paragraph{Magnitude pruning and the prune--retrain paradigm.}
\citet{han2015learning} established the template this paper works in: train a
dense network, remove the weights of smallest magnitude, and retrain the survivors
to recover the lost accuracy. \citet{zhu2018prune} replaced the one-shot removal
with a schedule, raising sparsity along a cubic ramp across training so that the
network is never asked to absorb a large discontinuity, and reported that the
interval between mask updates could vary over an order of magnitude with
negligible effect. \citet{gale2019state} audited that family at scale and
established the ramp-then-hold shape --- sparsity rises over the early portion of
training and the mask is then frozen while the weights recover --- that we adopt
in Section~\ref{sec:method}. Gradual magnitude pruning (GMP) in the sense of
those two papers is exactly what our \texttt{magnitude} criterion implements, a
correspondence we rely on in Section~\ref{sec:results} when we place our numbers
next to published ones.

\paragraph{Saliency criteria.}
SNIP \citep{lee2019snip} scores a weight by $|w \cdot \partial\mathcal{L}/\partial w|$,
the first-order effect of deleting it on the loss, and applies the score once
before training; SNIP-it \citep{verdenius2020snipit} applies the same criterion
iteratively during training, which is the form we use. WANDA
\citep{sun2024wanda} scores $|w_{ij}| \cdot \lVert X_j \rVert_2$, weighting
magnitude by the norm of the corresponding input activation over a calibration
set. WANDA's headline configuration ranks weights within each output row and its
authors argue that comparison group matters; that argument is made for large
language models, and their own appendix reports that the effect does not transfer
to image classifiers. We therefore use WANDA's \emph{metric} under the same
global threshold as the other two criteria, and present the grouping as our
design choice rather than on WANDA's authority.

\paragraph{Lottery tickets and data-free criteria.}
\citet{frankle2019lottery} showed that dense networks contain sparse subnetworks
that train to comparable accuracy from a rewound initialisation, which reframed
pruning as a search for structure rather than a compression post-process. GraSP
\citep{wang2020grasp} and SynFlow \citep{tanaka2020synflow} prune at
initialisation from gradient-flow and synaptic-flow signals respectively;
SynFlow additionally characterises layer collapse, the failure mode in which a
single layer is emptied and the network's function becomes constant, which is the
dominant risk at the sparsities we run. \citet{blalock2020state} document how
rarely pruning papers are mutually comparable --- differing protocols, missing
baselines, incommensurable sparsity conventions --- which is the direct motivation
for the matched-budget design in Section~\ref{sec:experiments}.

\paragraph{Dynamic sparse training.}
A parallel line keeps the parameter count fixed but lets the \emph{support} move.
SET \citep{mocanu2018set} prunes and regrows connections at random during
training; RigL \citep{evci2020rigl} regrows on gradient magnitude and introduces
the Erd\H{o}s--R\'{e}nyi-Kernel layer budget that most of this line now uses;
GraNet \citep{liu2021granet} interpolates between dense-to-sparse and
sparse-to-sparse training, and reports the strongest CIFAR-10 numbers we quote in
Section~\ref{sec:results}; EAST \citep{li2025east} targets the extreme regime
specifically and is, to our knowledge, the only published source reporting
CIFAR-10 accuracy for a ResNet at 99.9\% sparsity and beyond. This paper studies
\emph{static} pruning --- a mask that only ever removes weights --- so these
methods are context rather than baselines, and we do not claim results against
them. We list them because a reader evaluating an extreme-sparsity claim will
reasonably ask why the mask is not allowed to move, and the answer is scope: an
objective-level intervention and a support-level intervention are separable, and
we vary one of them.

\subsection{Contrastive learning}

SimCLR \citep{chen2020simclr} learns representations by pulling together two
augmented views of the same image and pushing apart views of different images
through the NT-Xent loss over in-batch negatives. MoCo \citep{he2020moco}
decouples the negative count from the batch by maintaining a momentum-updated
queue of keys. SupCon \citep{khosla2020supcon} generalises NT-Xent to multiple
positives per anchor by admitting every same-label sample into the positive set,
and shows that the formulation with the sum over positives outside the logarithm
is the better-behaved one; at a single positive it reduces to NT-Xent. These
three supply the loss forms that Section~\ref{sec:method} composes.

Contrastive objectives have also been used as distillation signals. CRD
\citep{tian2020crd} transfers a teacher's representation to a student by
maximising a lower bound on their mutual information, and is the closest
antecedent to using a contrastive term where a logit-matching term
\citep{hinton2015distilling} would ordinarily go.

\subsection{Contrastive objectives for compression}

CAP \citep{xu2022cap} is the direct antecedent of this work. It supervises a
model during gradual pruning with three contrastive terms, each contrasting the
current sparse student against a different frozen teacher --- the pretrained model
(PrC), the task-fine-tuned model (FiC), and previously saved sparse snapshots
(SnC) --- combined with cross-entropy into a single four-term weighted sum. CAP
reports on pre-trained language models (BERT-family encoders on GLUE tasks) at up
to 97\% sparsity.

\paragraph{Differences from CAP.}
The module decomposition used in this paper --- PrC, FiC and SnC, and the
four-term objective that combines them with cross-entropy --- is
\textbf{CAP's contribution, not ours}, and we do not present it as new. Our
objective is CAP's objective with the $\lambda$ indices permuted. The differences
between our instantiation and CAP's are the following, and none of them is a
claim of superiority in itself:

\begin{table}[t]
\centering
\small
\caption{Differences between CAP \citep{xu2022cap} and the instantiation studied
here. The three-module decomposition and the four-term objective are shared and
are CAP's.}
\label{tab:cap-diff}
\begin{tabular}{@{}p{0.16\textwidth}p{0.38\textwidth}p{0.38\textwidth}@{}}
\toprule
 & CAP \citep{xu2022cap} & This paper \\
\midrule
Modules & PrC, FiC, SnC & PrC, FiC, SnC --- identical \\
Objective & four-term weighted sum & same four terms, $\lambda$ indices permuted \\
Loss form & one functional, two positive sets, shared denominator & two functionals summed (SupCon $+$ NT-Xent), SimCLR-style all-pairs \\
Similarities & rectangular $B \times N$: student anchors $\times$ teacher candidates & square $2B \times 2B$ over $\mathrm{cat}[z_s, z_t]$ \\
Projection head & none; contrasts the \texttt{[CLS]} state & $\mathrm{Linear}(d,128)$, trainable, never pruned \\
Negatives & memory bank, $N = 4096$ pre-encoded & in-batch only, $2B$ \\
Modality & pre-trained language models & convolutional vision backbones \\
Max.\ sparsity reported & 97\% & 99.9\% \\
\bottomrule
\end{tabular}
\end{table}

Table~\ref{tab:cap-diff} sets them out. Two deserve comment. First, the
\emph{loss form}: CAP computes one functional over a rectangular
student-anchor $\times$ teacher-candidate similarity matrix, whereas the
instantiation we run stacks student and teacher embeddings and computes SupCon
and NT-Xent over the resulting square matrix, which admits within-student and
within-teacher blocks that CAP's formulation does not contain. This is a real
difference and not a reparameterisation; we ran both forms as a controlled
two-variable ablation and report the outcome in Section~\ref{sec:results}.
Second, the \emph{projection head}: CAP contrasts the backbone representation
directly, while we interpose a trainable linear head that is never pruned. This
is the substantive caveat on our claim and we return to it in
Section~\ref{sec:discussion}. CAP's memory bank is a component we omit rather
than an improvement we made; it slots into the rectangular formulation with no
change to the loss, and its absence reduces our negative count from $4096$ to
$2B$.

What remains, after those subtractions, is what we claim: the objective works on
vision backbones, works under three different pruning criteria, and continues to
work at sparsities beyond the range CAP examined.
